%----------------------------------------------------------------------------------------
% HYPOTHESETOETSEN. Alle opgaven zijn eigen werk.
%
% Opbouw: Start, leerstof (theorie, werkwijze of uitleg, voorbeeld, korte oefening), en op een
% nieuwe pagina de Opdrachten (door elkaar, stijgende moeilijkheid), Koppel en Test jezelf.
% Labels van dit hoofdstuk beginnen met h3-.
%----------------------------------------------------------------------------------------

\mychapter{3}{Hypothesetoetsen}

\lkr{Leerplan Basis: keuzedoel K5 (na te kijken) (hypothesen toetsen met nulhypothese, alternatieve hypothese, p-waarde en significantieniveau), met de wenken over de steekproevenverdeling, eenzijdige en tweezijdige toetsen, toetsen voor proporties en gemiddelden, de twee soorten fouten en de centrale limietstelling. Toetsen voor een proportie doen we exact met de binomiale verdeling; toetsen voor een gemiddelde met een bekende $\sigma$ en de normale verdeling. De t-toets valt buiten het leerplan. De p-waarde mag met ICT.}

Een producent beweert dat hoogstens 2~\% van zijn producten defect is. Je test er 200 en vindt er 9 defect. Liegt hij, of had je pech? Met een hypothesetoets beslis je dat op een vaste manier. Je gaat ervan uit dat de bewering klopt, en berekent hoe waarschijnlijk je resultaat dan is. Is dat heel onwaarschijnlijk, dan verwerp je de bewering. In dit hoofdstuk leer je die redenering, die in alle wetenschappen gebruikt wordt, en de valkuilen ervan.

\begin{oefeningen}[Start]
\oef{h3-start-1}

\oef{h3-verken}
\end{oefeningen}

\lkr{Laat oefening \ref{oef:h3-verken} maken vóór de uitleg. De redenering van vraag c (\textit{als de dobbelsteen eerlijk is, hoe groot is dan de kans op zoveel zessen?}) is de kern van elke toets.}

%========================================================================================
\section{De steekproevenverdeling}

In hoofdstuk 1 zag je dat een steekproefgrootheid schommelt van steekproef tot steekproef. Om te beoordelen of een resultaat verrassend is, moet je weten hoe sterk. Daarvoor bekijk je de steekproefgrootheid als toevalsvariabele, met een eigen kansverdeling.

\begin{theorie}
De \textbf{\textcolor{accent}{steekproevenverdeling}} van een steekproefgrootheid is haar kansverdeling over alle mogelijke aselecte steekproeven van grootte $n$.
\end{theorie}

\subsection{Een proportie}

Kies je aselect $n$ personen uit een grote populatie met aandeel $p$, dan heeft elke gekozen persoon (bijna) dezelfde kans $p$, onafhankelijk van de anderen. Het aantal $X$ met de eigenschap is dus binomiaal verdeeld, $X \sim B(n, p)$, met $\E(X) = np$ en $\Var(X) = np(1 - p)$. Het aandeel in de steekproef is $\hat{p} = \frac{X}{n}$.

\begin{theorie}
\begin{stelling}{Steekproevenverdeling van een proportie}
Voor aselecte steekproeven van grootte $n$ uit een grote populatie met proportie $p$ is

\vspace{1mm}
\qquad $\formule{\E(\hat{p}) = p}$ \qquad en \qquad $\formule{\sigma(\hat{p}) = \sqrt{\dfrac{p(1 - p)}{n}}}$
\end{stelling}

\begin{bewijs}
Neemt $X$ de waarde $k$ aan, dan neemt $\hat{p}$ de waarde $\frac{k}{n}$ aan, met dezelfde kans. Uit de definitie van de verwachtingswaarde volgt
\bewijsregel{\E(\hat{p}) = \sum_{k=0}^{n} \dfrac{k}{n}\, P(X = k) = \dfrac{1}{n} \sum_{k=0}^{n} k\, P(X = k) = \dfrac{1}{n}\, \E(X) = \dfrac{np}{n} = p .}
Op dezelfde manier is de variantie
\bewijsregel{\Var(\hat{p}) = \sum_{k=0}^{n} \Bigl(\dfrac{k}{n} - p\Bigr)^2 P(X = k) = \dfrac{1}{n^2} \sum_{k=0}^{n} (k - np)^2\, P(X = k) = \dfrac{\Var(X)}{n^2} ,}
en dat is $\frac{np(1 - p)}{n^2} = \frac{p(1 - p)}{n}$. De standaardafwijking is de vierkantswortel daarvan.
\end{bewijs}
\end{theorie}

De eerste formule zegt dat $\hat{p}$ gemiddeld juist zit: een aselecte steekproef schat $p$ zonder systematische fout. De tweede zegt hoe ver $\hat{p}$ typisch van $p$ ligt, en dat die afstand kleiner wordt als $n$ groeit. Is bovendien $np \geqslant 5$ en $n(1 - p) \geqslant 5$, dan is $X$, en dus ook $\hat{p}$, bij benadering normaal verdeeld.

\voorbeeld{}{In een stad doet 30~\% aan vrijwilligerswerk. Wat is de kans dat in een aselecte steekproef van 100 inwoners minstens 40~\% vrijwilligerswerk doet?\label{vb:h3-steekproef}}{%
\uitleg{We gebruiken de steekproevenverdeling van $\hat{p}$. Omdat $np = 30$ en $n(1 - p) = 70$ allebei minstens 5 zijn, is ze bij benadering normaal.}
\begin{stappen}
\sigma(\hat{p}) = \sqrt{\frac{0{,}3 \cdot 0{,}7}{100}} \approx 0{,}046 & Met $\E(\hat{p}) = 0{,}3$. \\[3ex]
P(\hat{p} \geqslant 0{,}4) \approx 0{,}015 & \texttt{normalcdf(0.4,1E99,0.3,0.0458)}: ongeveer 1,5~\%.
\end{stappen}
Een steekproef met 40~\% vrijwilligers is dus zeldzaam als $p = 0{,}3$. Vind je ze toch, dan twijfel je aan $p = 0{,}3$. Dat is het idee van een toets.}

\begin{oefeningen}
\oef{h3-phat}
\end{oefeningen}

\subsection{Een gemiddelde}

Meet je een getal bij elk lid, zoals een lengte of een wachttijd, dan heeft de populatie een gemiddelde $\mu$ en een standaardafwijking $\sigma$. Het gemiddelde $\overline{X}$ van een aselecte steekproef van $n$ waarden is een toevalsvariabele.

\begin{theorie}
\begin{stelling}{Steekproevenverdeling van een gemiddelde}
Voor aselecte steekproeven van grootte $n$ uit een grote populatie met gemiddelde $\mu$ en standaardafwijking $\sigma$ is

\vspace{1mm}
\qquad $\formule{\E(\overline{X}) = \mu}$ \qquad en \qquad $\formule{\sigma(\overline{X}) = \dfrac{\sigma}{\sqrt{n}}}$

\vspace{1mm}
Is de populatie normaal verdeeld, dan is ook $\overline{X}$ normaal verdeeld.
\end{stelling}
\end{theorie}

Het bewijs steunt op rekenregels voor sommen van onafhankelijke toevalsvariabelen, die buiten het leerplan vallen. In oefening \ref{oef:h3-dobbel} controleer je de formules voor twee worpen.

\naast{De stelling zegt dat steekproefgemiddelden rond $\mu$ liggen, en dichter naarmate $n$ groter is. Maar de spreiding daalt met $\sqrt{n}$, niet met $n$: vier keer zoveel waarnemingen halveert ze. De figuur toont de verdeling van $\overline{X}$ voor een normale populatie met $\mu = 70$ en $\sigma = 8$, bij steekproeven van 4, 16 en 64. De standaardafwijking is $\frac{8}{\sqrt n}$: 4, 2 en 1.}{\begin{tikzpicture}[x=.3cm,y=6cm,font=\scriptsize]
  \draw[->,grijs] (56,0) -- (84.5,0);
  \foreach \t in {60,70,80}{\node[below] at (\t,0) {\t};}
  \draw[grijs,thick,dashed] plot[domain=57:83,samples=80,smooth] (\x,{gauss(\x,70,4)});
  \draw[accent,thick] plot[domain=60:80,samples=80,smooth] (\x,{gauss(\x,70,2)});
  \draw[fel!80!black,thick] plot[domain=64:76,samples=80,smooth] (\x,{gauss(\x,70,1)});
  \node[fel!80!black,right] at (71.2,.38) {$n = 64$};
  \node[text=accent,right] at (72.5,.17) {$n = 16$};
  \node[right] at (75.5,.06) {$n = 4$};
\end{tikzpicture}}

\voorbeeld{}{De massa van pakken koffie is normaal verdeeld met $\mu = 250$~g en $\sigma = 4$~g. Een controleur weegt 16 aselect gekozen pakken. Wat is de kans dat hun gemiddelde onder 248~g ligt? En de kans dat één pak onder 248~g weegt?\label{vb:h3-xbar}}{%
\uitleg{We berekenen beide kansen, om het verschil te zien.}
\begin{stappen}
\sigma(\overline{X}) = \tfrac{4}{\sqrt{16}} = 1 & $\overline{X} \sim N(250, 1)$, want de populatie is normaal. \\[1ex]
P(\overline{X} < 248) \approx 0{,}023 & 248 ligt twee standaardafwijkingen onder 250. \\[1ex]
P(X < 248) \approx 0{,}31 & Voor één pak: \texttt{normalcdf(-1E99,248,250,4)}.
\end{stappen}
Eén licht pak is heel gewoon, een licht gemiddelde van 16 pakken niet.}

\begin{oefeningen}
\oef{h3-steekproef}

\oef{h3-dobbel}
\end{oefeningen}

\subsection{De centrale limietstelling}

Wachttijden, inkomens of bedragen op een kasticket zijn scheef verdeeld. De figuur toont een heel scheve populatie met $\mu = 1$, en de exacte verdeling van $\overline{X}$ bij $n = 4$ en $n = 16$: hoe groter $n$, hoe meer ze op een Gausskromme lijkt.

\figuren{\begin{tikzpicture}[x=1.35cm,y=.85cm,font=\scriptsize]
  \draw[->,grijs] (0,0) -- (3.2,0) node[right,text=black] {$x$};
  \draw[->,grijs] (0,0) -- (0,1.8);
  \foreach \t in {0,1,2,3}{\node[below] at (\t,0) {\t};}
  \draw[grijs,dashed] (1,0) -- (1,1.7);
  \draw[accent,thick] plot[smooth] coordinates {(0,1.000) (0.05,0.951) (0.1,0.905) (0.15,0.861) (0.2,0.819) (0.25,0.779) (0.3,0.741) (0.35,0.705) (0.4,0.670) (0.45,0.638) (0.5,0.607) (0.55,0.577) (0.6,0.549) (0.65,0.522) (0.7,0.497) (0.75,0.472) (0.8,0.449) (0.85,0.427) (0.9,0.407) (0.95,0.387) (1,0.368) (1.05,0.350) (1.1,0.333) (1.15,0.317) (1.2,0.301) (1.25,0.287) (1.3,0.273) (1.35,0.259) (1.4,0.247) (1.45,0.235) (1.5,0.223) (1.55,0.212) (1.6,0.202) (1.65,0.192) (1.7,0.183) (1.75,0.174) (1.8,0.165) (1.85,0.157) (1.9,0.150) (1.95,0.142) (2,0.135) (2.05,0.129) (2.1,0.122) (2.15,0.116) (2.2,0.111) (2.25,0.105) (2.3,0.100) (2.35,0.095) (2.4,0.091) (2.45,0.086) (2.5,0.082) (2.55,0.078) (2.6,0.074) (2.65,0.071) (2.7,0.067) (2.75,0.064) (2.8,0.061) (2.85,0.058) (2.9,0.055) (2.95,0.052) (3,0.050)};
  \node[anchor=north east] at (3.2,1.8) {populatie};
\end{tikzpicture}\quad
\begin{tikzpicture}[x=1.35cm,y=.85cm,font=\scriptsize]
  \draw[->,grijs] (0,0) -- (3.2,0) node[right,text=black] {$x$};
  \draw[->,grijs] (0,0) -- (0,1.8);
  \foreach \t in {0,1,2,3}{\node[below] at (\t,0) {\t};}
  \draw[grijs,dashed] (1,0) -- (1,1.7);
  \draw[accent,thick] plot[smooth] coordinates {(0,0.000) (0.05,0.004) (0.1,0.029) (0.15,0.079) (0.2,0.153) (0.25,0.245) (0.3,0.347) (0.35,0.451) (0.4,0.551) (0.45,0.643) (0.5,0.722) (0.55,0.787) (0.6,0.836) (0.65,0.870) (0.7,0.890) (0.75,0.896) (0.8,0.890) (0.85,0.874) (0.9,0.850) (0.95,0.818) (1,0.781) (1.05,0.741) (1.1,0.697) (1.15,0.652) (1.2,0.607) (1.25,0.561) (1.3,0.517) (1.35,0.474) (1.4,0.433) (1.45,0.394) (1.5,0.357) (1.55,0.322) (1.6,0.290) (1.65,0.261) (1.7,0.233) (1.75,0.209) (1.8,0.186) (1.85,0.165) (1.9,0.146) (1.95,0.130) (2,0.115) (2.05,0.101) (2.1,0.089) (2.15,0.078) (2.2,0.068) (2.25,0.060) (2.3,0.052) (2.35,0.046) (2.4,0.040) (2.45,0.035) (2.5,0.030) (2.55,0.026) (2.6,0.023) (2.65,0.020) (2.7,0.017) (2.75,0.015) (2.8,0.013) (2.85,0.011) (2.9,0.010) (2.95,0.008) (3,0.007)};
  \node[anchor=north east] at (3.2,1.8) {$\overline{X}$, $n = 4$};
\end{tikzpicture}\quad
\begin{tikzpicture}[x=1.35cm,y=.85cm,font=\scriptsize]
  \draw[->,grijs] (0,0) -- (3.2,0) node[right,text=black] {$x$};
  \draw[->,grijs] (0,0) -- (0,1.8);
  \foreach \t in {0,1,2,3}{\node[below] at (\t,0) {\t};}
  \draw[grijs,dashed] (1,0) -- (1,1.7);
  \draw[accent,thick] plot[smooth] coordinates {(0.3,0.002) (0.35,0.008) (0.4,0.025) (0.45,0.066) (0.5,0.144) (0.55,0.271) (0.6,0.449) (0.65,0.671) (0.7,0.916) (0.75,1.158) (0.8,1.370) (0.85,1.529) (0.9,1.619) (0.95,1.637) (1,1.587) (1.05,1.483) (1.1,1.339) (1.15,1.172) (1.2,0.997) (1.25,0.826) (1.3,0.669) (1.35,0.529) (1.4,0.410) (1.45,0.312) (1.5,0.233) (1.55,0.171) (1.6,0.124) (1.65,0.088) (1.7,0.062) (1.75,0.043) (1.8,0.030) (1.85,0.020) (1.9,0.013) (1.95,0.009) (2,0.006) (2.05,0.004) (2.1,0.002)};
  \node[anchor=north east] at (3.2,1.8) {$\overline{X}$, $n = 16$};
\end{tikzpicture}}

\begin{theorie}
\begin{stelling}{Centrale limietstelling}
Voor grote $n$ is $\overline{X}$ bij benadering normaal verdeeld, met gemiddelde $\mu$ en standaardafwijking $\frac{\sigma}{\sqrt{n}}$, welke verdeling de populatie ook heeft:

\vspace{1mm}
\qquad $\formule{\overline{X} \approx N\Bigl(\mu, \dfrac{\sigma}{\sqrt{n}}\Bigr)}$
\end{stelling}
\end{theorie}

Het bewijs gebruikt wiskunde van na het secundair onderwijs. Als vuistregel volstaat $n \geqslant 30$, tenzij de populatie heel scheef is.

\voorbeeld{}{Het bedrag op een kasticket in een supermarkt is scheef verdeeld, met $\mu = 38$~euro en $\sigma = 24$~euro. Wat is de kans dat 64 aselect gekozen kastickets gemiddeld meer dan 42~euro bedragen?\label{vb:h3-clt}}{%
\uitleg{De populatie is scheef, maar $n = 64$ is groot: dus is $\overline{X}$ bij benadering normaal.}
\begin{stappen}
\sigma(\overline{X}) = \tfrac{24}{\sqrt{64}} = 3 & $\overline{X} \approx N(38, 3)$. \\[1ex]
P(\overline{X} > 42) \approx 0{,}091 & \texttt{normalcdf(42,1E99,38,3)}.
\end{stappen}
Voor \textit{één} kasticket lukt dat niet: één bedrag is niet normaal verdeeld.}

\begin{oefeningen}
\oef{h3-clt}
\end{oefeningen}

%========================================================================================
\section{De toets}

\subsection{Nulhypothese en alternatieve hypothese}

In de Verken beweert Sam dat de dobbelsteen vervalst is. Om dat te onderzoeken, zet je twee uitspraken tegenover elkaar: de dobbelsteen is eerlijk, of hij gooit te vaak zes. Een toets kiest tussen die twee op basis van de data.

\begin{theorie}
De \textbf{\textcolor{accent}{nulhypothese}} $\formule{H_0}$ is de bewering die je toetst, met een vaste waarde van de parameter, zoals $p = \frac16$. De \textbf{\textcolor{accent}{alternatieve hypothese}} $\formule{H_1}$ is wat je vermoedt als $H_0$ niet klopt, zoals $p > \frac16$.
\end{theorie}

$H_0$ is de saaie uitspraak: er is geen effect, de bewering klopt. $H_1$ is wat de onderzoeker wil aantonen. In $H_0$ staat een gelijkheidsteken: onder $H_0$ moet je kansen kunnen berekenen.

\voorbeeld{}{Een fabrikant beweert dat 90~\% van zijn zonnepanelen na 20 jaar nog werkt. Een consumentenorganisatie vermoedt dat het minder is. Formuleer $H_0$ en $H_1$.\label{vb:h3-hypothesen}}{%
\uitleg{We zoeken eerst de parameter, daarna de bewering met een vaste waarde, en tot slot het vermoeden.}
\begin{stappenbreed}
p = \text{aandeel van alle panelen dat na 20 jaar werkt} & De parameter waarover de discussie gaat. \\[.6ex]
H_0\colon p = 0{,}9 & De bewering van de fabrikant. \\[.6ex]
H_1\colon p < 0{,}9 & Het vermoeden van de organisatie.
\end{stappenbreed}}

\begin{oefeningen}
\oef{h3-hypothesen}
\end{oefeningen}

\subsection{De p-waarde en het besluit}

Een toets is een bewijs uit het ongerijmde met kansen. Je neemt aan dat $H_0$ klopt, en berekent hoe waarschijnlijk een resultaat als het jouwe dan is. Is dat heel onwaarschijnlijk, dan verwerp je de aanname.

\begin{theorie}
De \textbf{\textcolor{accent}{p-waarde}} is de kans, als $H_0$ waar is, op een steekproefresultaat dat minstens zo extreem is als het waargenomen resultaat (in de richting van $H_1$).

\vspace{1mm}
Het \textbf{\textcolor{accent}{significantieniveau}} $\formule{\alpha}$ kies je vooraf, meestal $0{,}05$. Is de p-waarde $\leqslant \alpha$, dan \textbf{\textcolor{accent}{verwerp}} je $H_0$: het resultaat is \textbf{\textcolor{accent}{significant}}. Anders verwerp je $H_0$ niet.
\end{theorie}

\voorbeeld{}{In de Verken gooit een dobbelsteen 16 keer zes bij 60 worpen. Bij $H_0\colon p = \frac16$ en $H_1\colon p > \frac16$ is de p-waarde $P(X \geqslant 16) \approx 0{,}034$. Wat besluit je bij $\alpha = 0{,}05$ en bij $\alpha = 0{,}01$? Wat betekent die p-waarde?\label{vb:h3-pwaarde}}{%
\uitleg{We vergelijken de p-waarde met elk niveau, en zeggen in de context wat ze meet.}
\begin{stappenbreed}
0{,}034 \leqslant 0{,}05\colon\ H_0 \text{ verwerpen} & Bij $\alpha = 0{,}05$ is het resultaat significant. \\[1ex]
0{,}034 > 0{,}01\colon\ H_0 \text{ niet verwerpen} & Bij het strengere niveau volstaat het niet.
\end{stappenbreed}
Betekenis: als de dobbelsteen eerlijk is, gooit hij in ongeveer 3,4~\% van de reeksen van 60 worpen minstens 16 keer zes.}

Twee klassieke denkfouten. Verwerp je $H_0$ niet, dan is ze niet bewezen: er zijn alleen te weinig aanwijzingen tegen haar. En de p-waarde is niet de kans dat $H_0$ waar is: ze is net berekend in de veronderstelling dat de dobbelsteen eerlijk is.

\begin{oefeningen}
\oef{h3-begrippen}
\end{oefeningen}

\subsection{Een toets voor een proportie}

Bij een proportie is het aantal successen onder $H_0$ binomiaal verdeeld: de p-waarde berekenen we exact.

\voorbeeld{}{Een speler vermoedt dat een dobbelsteen te vaak zes gooit. Bij 120 worpen valt 30 keer zes. Toets bij $\alpha = 0{,}05$.\label{vb:h3-proportie}}{%
\uitleg{Bij een eerlijke dobbelsteen is $p = \frac16$, en verwacht je 20 zessen. We berekenen hoe waarschijnlijk 30 of meer zessen zijn als de dobbelsteen eerlijk is.}
\begin{stappenbreed}
H_0\colon p = \tfrac16 \qquad H_1\colon p > \tfrac16 & De speler vermoedt te veel zessen. \\[1ex]
\text{onder } H_0\colon \ X \sim B(120,\, \tfrac16) & $X$ is het aantal zessen. \\[1ex]
\text{p-waarde} = P(X \geqslant 30) = 1 - P(X \leqslant 29) \approx 0{,}013 & Met \texttt{binomcdf(120,1/6,29)}. \\[1ex]
0{,}013 \leqslant 0{,}05: \ H_0 \text{ verwerpen} & Significant: de dobbelsteen gooit waarschijnlijk te vaak zes.
\end{stappenbreed}}

\begin{werkwijze}
\textbf{1.}~Formuleer $H_0$ (met $=$) en $H_1$ ($>$, $<$ of $\neq$), en kies $\alpha$. \ \textbf{2.}~Bepaal de verdeling van de toetsingsgrootheid als $H_0$ waar is. \ \textbf{3.}~Bereken de p-waarde: de kans op het waargenomen resultaat of iets extremer, in de richting van $H_1$. \ \textbf{4.}~Vergelijk met $\alpha$. \ \textbf{5.}~Formuleer het besluit in de context, voorzichtig.
\end{werkwijze}

\begin{oefeningen}
\oef{h3-proportie}
\end{oefeningen}

\subsection{Een toets voor een gemiddelde}

Bij een gemiddelde gebruik je de steekproevenverdeling van $\overline{X}$. In dit hoofdstuk kennen we de standaardafwijking $\sigma$ van de populatie, bijvoorbeeld uit vroegere metingen.

\voorbeeld{}{Een machine vult flessen met gemiddeld 500~ml, met $\sigma = 2$~ml. Een consumentenorganisatie meet 25 flessen en vindt een gemiddelde van 499,1~ml. Vult de machine gemiddeld te weinig? Toets bij $\alpha = 0{,}05$.\label{vb:h3-gemiddelde}}{%
\uitleg{Als de machine goed werkt, is $\overline{X}$ normaal verdeeld met gemiddelde 500 en standaardafwijking $\frac{2}{\sqrt{25}} = 0{,}4$. We berekenen de kans op een gemiddelde van 499,1 of minder.}
\begin{stappenbreed}
H_0\colon \mu = 500 \qquad H_1\colon \mu < 500 & Het vermoeden: te weinig. \\[1ex]
\text{onder } H_0\colon \ \overline{X} \sim N(500;\, 0{,}4) & $\sigma(\overline{X}) = \frac{2}{\sqrt{25}}$. \\[1ex]
\text{p-waarde} = P(\overline{X} \leqslant 499{,}1) \approx 0{,}012 & $z = \frac{499{,}1 - 500}{0{,}4} = -2{,}25$. \\[1ex]
0{,}012 \leqslant 0{,}05: \ H_0 \text{ verwerpen} & De machine vult waarschijnlijk gemiddeld te weinig.
\end{stappenbreed}}

\nujij{Bij een toets is de p-waarde 0,08. Wat besluit je bij $\alpha = 0{,}05$? Betekent dit dat $H_0$ waar is?}{$0{,}08 > 0{,}05$: $H_0$ niet verwerpen. Dat bewijst niet dat $H_0$ waar is; de aanwijzingen ertegen zijn alleen niet sterk genoeg.}

\begin{oefeningen}
\oef{h3-toets}
\end{oefeningen}

%========================================================================================
\section{Eenzijdig en tweezijdig}

Vermoed je alleen dat een parameter \textit{anders} is, zonder richting, dan telt een extreem resultaat aan beide kanten even zwaar. Dan moet de p-waarde ook beide kanten meetellen.

\begin{theorie}
Een toets is \textbf{\textcolor{accent}{eenzijdig}} als $H_1$ van de vorm $p > p_0$ of $p < p_0$ is, en \textbf{\textcolor{accent}{tweezijdig}} als $H_1$ van de vorm $\formule{p \neq p_0}$ is (en zo ook voor $\mu$).

\vspace{1mm}
Bij een tweezijdige toets is

\vspace{1mm}
\qquad $\formule{\text{p-waarde} = 2 \cdot \text{de kans in de staart aan de kant van het resultaat}}$
\end{theorie}

Waarom het dubbel? Een resultaat dat even ver aan de andere kant ligt, is even sterk bewijs tegen $H_0$. Bij een symmetrische verdeling, zoals de normale of $B(n;\, 0{,}5)$, heeft die andere staart precies dezelfde kans. Bij een binomiale verdeling met $p_0 \neq 0{,}5$ is de verdeling niet symmetrisch: daar is het verdubbelen een afspraak.

\voorbeeld{}{Een munt landt 60 keer op kop bij 100 worpen. Is de munt oneerlijk? Toets tweezijdig bij $\alpha = 0{,}05$.\label{vb:h3-tweezijdig}}{%
\uitleg{\textit{Oneerlijk} kan naar kop of naar munt: $H_1\colon p \neq 0{,}5$. Het resultaat ligt aan de kant van te veel kop; we verdubbelen die staart.}
\begin{stappenbreed}
H_0\colon p = 0{,}5 \qquad H_1\colon p \neq 0{,}5 & Tweezijdig. \\[1ex]
P(X \geqslant 60) \approx 0{,}0284 \quad \Rightarrow \quad \text{p-waarde} \approx 2 \cdot 0{,}0284 \approx 0{,}057 & Met $X \sim B(100;\, 0{,}5)$; 40 of minder keer kop is even extreem. \\[1ex]
0{,}057 > 0{,}05: \ H_0 \text{ niet verwerpen} & Niet significant.
\end{stappenbreed}}

Merk op: had je vooraf vermoed dat de munt te vaak kop geeft, dan was de p-waarde 0,028 en had je $H_0$ verworpen. Daarom kies je $H_1$ \textit{voor} je de data bekijkt, op basis van de vraag, en niet op basis van het resultaat.

\begin{oefeningen}
\oef{h3-tweezijdig}

\oef{h3-tweezijdig-2}
\end{oefeningen}

%========================================================================================
\section{Twee soorten fouten}

\subsection{Fouten van de eerste en de tweede soort}

Een toets beslist op basis van een steekproef, en kan zich dus vergissen. Er zijn twee manieren.

\begin{theorie}
{\centering\arrayrulecolor{accent}\renewcommand{\arraystretch}{1.6}%
\begin{tabular}{|>{\centering\arraybackslash}m{3.6cm}|>{\centering\arraybackslash}m{4.4cm}|>{\centering\arraybackslash}m{4.4cm}|}
\hline
 & \textbf{$H_0$ is waar} & \textbf{$H_0$ is niet waar} \\ \hline
\textbf{$H_0$ verwerpen} & \textbf{\textcolor{accent}{fout van de eerste soort}} (type I) & juist \\ \hline
\textbf{$H_0$ niet verwerpen} & juist & \textbf{\textcolor{accent}{fout van de tweede soort}} (type II) \\ \hline
\end{tabular}\par}
\end{theorie}

Bij een medicijn dat getest wordt ($H_0$: het werkt niet) is een fout van de eerste soort: een medicijn goedkeuren dat niet werkt. Een fout van de tweede soort: een werkzaam medicijn afwijzen.

\begin{oefeningen}
\oef{h3-fouten}
\end{oefeningen}

\subsection{De kans op een fout van de eerste soort}

Het significantieniveau is meer dan een grens voor de p-waarde: het is ook het risico op een vals alarm dat je vooraf aanvaardt.

\begin{theorie}
\begin{stelling}{Fout van de eerste soort}
Verwerp je $H_0$ als de p-waarde $\leqslant \alpha$, dan is de kans op een fout van de eerste soort hoogstens $\formule{\alpha}$.
\end{stelling}

\begin{bewijs}
Noem $T$ de toetsingsgrootheid, zoals het aantal successen of $\overline{X}$, en stel dat $H_0$ waar is. Neem eerst $H_1$ met $>$. Bij een uitkomst $t$ is de p-waarde $P(T \geqslant t)$, en die kans wordt kleiner als $t$ groter wordt. Is $c$ de kleinste uitkomst met $P(T \geqslant c) \leqslant \alpha$, dan is de p-waarde dus $\leqslant \alpha$ precies als $T \geqslant c$ (bestaat zo'n $c$ niet, dan verwerp je nooit), en
\bewijsregel{P(H_0 \text{ verwerpen}) = P(T \geqslant c) \leqslant \alpha .}
Bij $H_1$ met $<$ is alles gespiegeld: je verwerpt precies als $T \leqslant c'$, met $P(T \leqslant c') \leqslant \alpha$. Bij $H_1$ met $\neq$ verwerp je als de staartkans aan de kant van $t$ hoogstens $\frac{\alpha}{2}$ is, dus als $T \geqslant c$ met $P(T \geqslant c) \leqslant \frac{\alpha}{2}$ of als $T \leqslant c'$ met $P(T \leqslant c') \leqslant \frac{\alpha}{2}$. Met de algemene somregel is
\bewijsregel{P(H_0 \text{ verwerpen}) \leqslant P(T \geqslant c) + P(T \leqslant c') \leqslant \dfrac{\alpha}{2} + \dfrac{\alpha}{2} = \alpha .}
\end{bewijs}
\end{theorie}

De uitkomsten waarbij je $H_0$ verwerpt, zoals $T \geqslant c$, vormen het \textbf{verwerpingsgebied}. Bij een gemiddelde is de kans erop precies $\alpha$; bij een aantal kan ze kleiner zijn, omdat een aantal met sprongen gaat.

\voorbeeld{}{Bij 20 worpen met een munt verwerp je $H_0\colon p = 0{,}5$ ten voordele van $H_1\colon p > 0{,}5$ als er minstens 15 keer kop valt. Wat is de kans op een fout van de eerste soort? En op een fout van de tweede soort als in werkelijkheid $p = 0{,}7$?\label{vb:h3-fouten}}{%
\uitleg{Een fout van de eerste soort: $H_0$ is waar, en toch valt 15 of meer keer kop. Een fout van de tweede soort: $p = 0{,}7$, en toch valt minder dan 15 keer kop.}
\begin{stappenbreed}
P(X \geqslant 15 \mid p = 0{,}5) \approx 0{,}021 & De kans op een fout van de eerste soort: kleiner dan 0,05. \\[1ex]
P(X \leqslant 14 \mid p = 0{,}7) \approx 0{,}58 & Bij $p = 0{,}7$ mis je de afwijking in 58~\% van de gevallen.
\end{stappenbreed}
Met maar 20 worpen ontdek je een munt die 70~\% kop geeft dus maar in 42~\% van de gevallen.}

\subsection{De fout van de tweede soort verkleinen}

Kies je $\alpha$ kleiner, dan verwerp je minder snel ten onrechte, maar schuift het verwerpingsgebied op en mis je vaker een echt effect: de kans op een fout van de tweede soort stijgt. Beide fouten tegelijk verkleinen kan alleen met een grotere steekproef.

\voorbeeld{}{Je toetst dezelfde munt nu met 50 worpen, opnieuw met $H_1\colon p > 0{,}5$ en $\alpha = 0{,}05$. Bepaal het verwerpingsgebied, en de kans op een fout van de tweede soort als $p = 0{,}7$.\label{vb:h3-fouten-n}}{%
\uitleg{We zoeken het kleinste aantal $c$ met $P(X \geqslant c) \leqslant 0{,}05$ onder $H_0$, en berekenen daarna de kans om onder $c$ te blijven als $p = 0{,}7$.}
\begin{stappenbreed}
P(X \geqslant 31) \approx 0{,}059, \quad P(X \geqslant 32) \approx 0{,}032 & Met $X \sim B(50;\, 0{,}5)$: verwerpen vanaf 32 keer kop. \\[1ex]
P(X \leqslant 31 \mid p = 0{,}7) \approx 0{,}14 & \texttt{binomcdf(50,0.7,31)}.
\end{stappenbreed}
In beide gevallen is de kans op een vals alarm hoogstens 5~\%, maar de kans om de oneerlijke munt te missen daalt van 58~\% naar 14~\%.}

\begin{oefeningen}
\oef{h3-fouten-2}

\oef{h3-fouten-n}
\end{oefeningen}

%========================================================================================
\opdrachtenkop

\begin{oefeningen}
\needspace{18\baselineskip}
\oef{h3-o3}

\needspace{12\baselineskip}
\oef{h3-d2}

\oef{h3-o1}

\needspace{16\baselineskip}
\oef{h3-o12}

\needspace{8\baselineskip}
\oef{h3-d3}

\oef{h3-o2}

\needspace{12\baselineskip}
\oef{h3-o11}

\needspace{10\baselineskip}
\oef{h3-o10}

\oef{h3-o13}

\oef{h3-o4}

\needspace{16\baselineskip}
\oef{h3-o14}

\needspace{12\baselineskip}
\oef{h3-o18}

\needspace{12\baselineskip}
\oef{h3-o16}

\needspace{10\baselineskip}
\oef{h3-o6}

\needspace{12\baselineskip}
\oef{h3-o15}

\needspace{12\baselineskip}
\oef{h3-o19}

\oef{h3-o7}

\needspace{12\baselineskip}
\oef{h3-o20}

\oef{h3-o21}

\needspace{12\baselineskip}
\oef{h3-o17}

\oef{h3-o5}

\oef{h3-o9}

\needspace{12\baselineskip}
\oef{h3-o23}

\oef{h3-o8}

\needspace{14\baselineskip}
\oef{h3-o22}

\needspace{12\baselineskip}
\oef{h3-o24}
\end{oefeningen}

%----------------------------------------------------------------------------------------
% KOPPEL: verbanden met andere leerstof
%----------------------------------------------------------------------------------------
\needspace{20\baselineskip}
\opdrachtsectie{Koppel}

\lkr{De oefening gebruikt voorwaardelijke kansen uit Kansrekenen.}

\begin{oefeningen}

\oef{h3-koppel-2}
\end{oefeningen}

%----------------------------------------------------------------------------------------
\opdrachtsectie{Test jezelf}

\begin{oefeningen}
\oef{h3-t1}

\oef{h3-t2}

\oef{h3-t3}

\oef{h3-t4}

\oef{h3-t5}

\oef{h3-t6}

\oef{h3-t7}
\end{oefeningen}

\herhaalschema{%
  \fv{oef:h3-t1}{vb:h3-proportie}{oef:h3-proportie}%
  \fv{oef:h3-t2}{vb:h3-gemiddelde}{oef:h3-toets}%
  \fv{oef:h3-t3}{vb:h3-tweezijdig}{oef:h3-tweezijdig-2}%
  \fzv{oef:h3-t4}{oef:h3-fouten}%
  \fv{oef:h3-t5}{vb:h3-xbar}{oef:h3-steekproef}%
  \fv{oef:h3-t6}{vb:h3-pwaarde}{oef:h3-begrippen}%
  \fv{oef:h3-t7}{vb:h3-fouten}{oef:h3-fouten-2}}

\ifzelfstudie
\section*{Maak je eigen samenvatting}
Sluit dit boek en schrijf in je schrift: de steekproevenverdeling van een proportie en van een gemiddelde, de centrale limietstelling, de begrippen nulhypothese, alternatieve hypothese, p-waarde en significantieniveau, de vijf stappen van een toets, het verschil tussen eenzijdig en tweezijdig, en de twee soorten fouten. Schrijf ook op wat het betekent als je $H_0$ niet verwerpt. Vergelijk daarna met de kaders.
\fi
